
\documentclass{article}
\usepackage{colortbl}
\usepackage{makecell}
\usepackage{multirow}
\usepackage{supertabular}

\begin{document}

\newcounter{utterance}

\centering \large Interaction Transcript for game `cladder', experiment `full\_v1.5\_default', episode 9292 with refineonly{-}Qwen3.5{-}27B.
\vspace{24pt}

{ \footnotesize  \setcounter{utterance}{1}
\setlength{\tabcolsep}{0pt}
\begin{supertabular}{c@{$\;$}|p{.15\linewidth}@{}p{.15\linewidth}p{.15\linewidth}p{.15\linewidth}p{.15\linewidth}p{.15\linewidth}}
    \# & \multicolumn{2}{c}{Player} && \multicolumn{2}{c}{Game Master} \\
    \hline

    \theutterance \stepcounter{utterance}  
    & & & \multicolumn{4}{p{0.6\linewidth}}{
        \cellcolor[rgb]{0.9,0.9,0.9}{
            \makecell[{{p{\linewidth}}}]{
                \texttt{\tiny{[P1$\langle$GM]}}
                \texttt{You are an expert in causal inference. The following question is not a typical commonsense query, but rather a meticulously designed question created by a professor specializing in causal inference, intended to assess the students' mastery of the course content.} \\
\texttt{Imagine a self{-}contained, hypothetical world with only the following conditions, and without any unmentioned factors or causal relationships: Zory has a direct effect on zuph and glimx. Zuph has a direct effect on jyka. Jyka has a direct effect on glimx. Zory is unobserved. The overall probability of zuph is 36\%. The probability of not zuph and glimx is 30\%. The probability of zuph and glimx is 31\%. Is the chance of glimx smaller when observing zuph?} \\
\texttt{Start your answer with 'Yes' or 'No', followed by additional reasoning or evidence to support your explanation.} \\
            }
        }
    }
     \\ \\

    \theutterance \stepcounter{utterance}  
    & \multicolumn{4}{p{0.6\linewidth}}{
        \cellcolor[rgb]{0.561,0.737,0.561}{
            \makecell[{{p{\linewidth}}}]{
                \texttt{\tiny{[P1$\rangle$GM]}}
                \texttt{No} \\
\\ 
\texttt{To determine whether the chance of *glimx* is smaller when observing *zuph*, we must calculate and compare the conditional probabilities \$P(\textbackslash{}text\{glimx\} \textbackslash{}mid \textbackslash{}text\{zuph\})\$ and \$P(\textbackslash{}text\{glimx\} \textbackslash{}mid \textbackslash{}neg\textbackslash{}text\{zuph\})\$ using the provided data.} \\
\\ 
\texttt{First, we identify the given probabilities:} \\
\texttt{{-} \$P(\textbackslash{}text\{zuph\}) = 36\textbackslash{}\%\$} \\
\texttt{{-} \$P(\textbackslash{}neg\textbackslash{}text\{zuph\} \textbackslash{}cap \textbackslash{}text\{glimx\}) = 30\textbackslash{}\%\$} \\
\texttt{{-} \$P(\textbackslash{}text\{zuph\} \textbackslash{}cap \textbackslash{}text\{glimx\}) = 31\textbackslash{}\%\$} \\
\\ 
\texttt{Next, we calculate the conditional probability of *glimx* given *zuph*:} \\
\texttt{\$\$P(\textbackslash{}text\{glimx\} \textbackslash{}mid \textbackslash{}text\{zuph\}) = \textbackslash{}frac\{P(\textbackslash{}text\{zuph\} \textbackslash{}cap \textbackslash{}text\{glimx\})\}\{P(\textbackslash{}text\{zuph\})\} = \textbackslash{}frac\{31\textbackslash{}\%\}\{36\textbackslash{}\%\} \textbackslash{}approx 0.861 \textbackslash{}text\{ (or \} 86.1\textbackslash{}\%)\$\$} \\
\\ 
\texttt{Then, we calculate the conditional probability of *glimx* given *not zuph* (\$\textbackslash{}neg\textbackslash{}text\{zuph\}\$). First, we determine \$P(\textbackslash{}neg\textbackslash{}text\{zuph\})\$:} \\
\texttt{\$\$P(\textbackslash{}neg\textbackslash{}text\{zuph\}) = 1 {-} P(\textbackslash{}text\{zuph\}) = 1 {-} 36\textbackslash{}\% = 64\textbackslash{}\%\$\$} \\
\texttt{Now, we compute the conditional probability:} \\
\texttt{\$\$P(\textbackslash{}text\{glimx\} \textbackslash{}mid \textbackslash{}neg\textbackslash{}text\{zuph\}) = \textbackslash{}frac\{P(\textbackslash{}neg\textbackslash{}text\{zuph\} \textbackslash{}cap \textbackslash{}text\{glimx\})\}\{P(\textbackslash{}neg\textbackslash{}text\{zuph\})\} = \textbackslash{}frac\{30\textbackslash{}\%\}\{64\textbackslash{}\%\} \textbackslash{}approx 0.469 \textbackslash{}text\{ (or \} 46.9\textbackslash{}\%)\$\$} \\
\\ 
\texttt{Comparing the two results, \$P(\textbackslash{}text\{glimx\} \textbackslash{}mid \textbackslash{}text\{zuph\}) \textbackslash{}approx 86.1\textbackslash{}\%\$ is significantly larger than \$P(\textbackslash{}text\{glimx\} \textbackslash{}mid \textbackslash{}neg\textbackslash{}text\{zuph\}) \textbackslash{}approx 46.9\textbackslash{}\%\$. Therefore, the chance of *glimx* is not smaller when observing *zuph*; it is actually larger. The causal structure involving the unobserved variable *Zory* creates a confounding effect (or selection bias depending on the specific graph interpretation), but the direct calculation of the conditional probabilities from the provided marginals definitively answers the question.} \\
            }
        }
    }
    & & \\ \\

    \theutterance \stepcounter{utterance}  
    & & & \multicolumn{2}{p{0.3\linewidth}}{
        \cellcolor[rgb]{0.9,0.9,0.9}{
            \makecell[{{p{\linewidth}}}]{
                \texttt{\tiny{[GM$|$GM]}}
                \texttt{No} \\
            }
        }
    }
    & & \\ \\

    \theutterance \stepcounter{utterance}  
    & & & \multicolumn{2}{p{0.3\linewidth}}{
        \cellcolor[rgb]{0.9,0.9,0.9}{
            \makecell[{{p{\linewidth}}}]{
                \texttt{\tiny{[GM$|$GM]}}
                \texttt{no} \\
            }
        }
    }
    & & \\ \\

    \theutterance \stepcounter{utterance}  
    & & & \multicolumn{2}{p{0.3\linewidth}}{
        \cellcolor[rgb]{0.9,0.9,0.9}{
            \makecell[{{p{\linewidth}}}]{
                \texttt{\tiny{[GM$|$GM]}}
                \texttt{game\_result = WIN} \\
            }
        }
    }
    & & \\ \\

\end{supertabular}
}

\end{document}
